\documentclass[english,10pt,a4paper]{article}
\usepackage[T1]{fontenc}
\usepackage{graphicx}
\usepackage{xcolor}
\usepackage{babel}
\title{Rulebook}

\author{Alexander Mödder and Alexander Vinkeloe}

%Design structure of this documnt:
%Show what exists first
%Explain how it is used second
%Go top down (Complex overarching to simple detail)
%? Try to not crossref? People can remember 3 things at most, if you have to make them

\begin{document}
	\maketitle
	
	%\section{Foreword} This one wont have a section formatting, but still needs to be there
	Game Overview\\
	
	Wheel of Elements is a Slay-the-Spire-like cooperative multi-classing card game. Become an indomitable tank, a Mage of Immense arcane Power or something in between. Beat four battles of increasing difficulty and prove your mettle.\\
	Wheel of Elements is a multi-classing game. Each player starts with two classes and a race. Improve or add a class after each won combat.\\
	Wheel of Elements is a card game. Each Class an Race gives you cards which combine to a 15 card Deck. Choose which you want to play.\\ 
	
	Wheel of Elements has an elemental Wheel. Create elements to consume them for Bursts of Power.\\
	
	Playtime\\
	
	The Game consists of four fights with deckbuilding sessions between them. A full game takes about 2 hours. Plan for some extra time your first game.\\
	
	Winning \& Losing\\
	
	Victory - Victory is achieved when each enemy is killed and at least one player survived.\\
	
	Loss - You lose when each player died in the same fight.\\
	
	\tableofcontents
	\section{Components}
	%recalc way at the end
	1 Rulebook
	1 Strategy Guide ??????????
	1 Elemental Wheel
	25 Element Chips
	10 Race Tableau
	30 Class Shingles
	5 Player Life Counter (Could be on the Race Tableau)
	~220 Small Cards 	(5 Base modifier decks a 10 cards)
						(10 Race modifier Decks a ~15 cards)
						(~20 From Passives)
	~210 Player Cards	(30 Class decks a 5 cards)
						(10 Race decks a 6 cards)
	20 Events
	150 Monsters
	5 Dice (Monster Die and Drunkenmaster Die)
	20 Plastic Cubes 	(Damage Trackers for monsters)
	50 Tokens Boons\&Banes
	
	
	\section{Setup}
	1.Pick a Race
	2.Pick a Level 1 Class
	3.Pick a different level 1 Class
	4.Shuffle your Action decks together
	5.Do the modifier changes to your modifier Deck depending of the Race Tableau and shuffle it. Then put the class shingles in the Race Tableau
	6.Put the Elemental Wheel in the middle
	7.Each player draws one Fight 1 encounter card
	8.Each player setups the his enemies corresponding to his Fight 1 encounter card
	9. Start Combat
	
	\section{Play Area}
	
	Action draw pile - Shuffle your deck before combat begins to create your draw pile
	Action discard pile - After you play a card, it goes here
	Modifier draw pile - Shuffle your deck before combat begins to create your draw pile
	Modifier discard pile - After you modify an attack, all drawn modifiers go here
	Tableau - The tableau tracks you level up decision and life
	Enemy Zone - Enemies go here, with the configuration described in the storybook
	
	
	\section{Your Decks}
	\subsection{The action Deck}
	Card pool - The cards that can be added to your deck
	Before combat select 15 cards, these will be your deck
	The action deck contains all cards you can play during combat.
	Card rarity - Cards can have different rarity levels, from bronze to mithril.
	Card power varies greatly, but two cards of the same rarity are generally equal
	
	\subsection{The modifier Deck}
	The modifier cards are used every time you attack an enemy. They can have three effects:
	Numerical - Adds a number to your attack damage
	Elements - Create the shown elements
	Draw - Perform all other effects of this modifier, and draw another
	
	
	\section{Leveling Up}
	
	When you level up, you may choose a new class and add its shingle to your tableau.
	The class shingles have to be placed in the leftmost slot of their rank.
	Rank 1 classes go in the top row.
	Rank 2 classes go in the middle row.
	Rank 3 classes go in the bottom. %Picture for this
	Each slot for a shingle has changes to the modifier deck printed on it.
	When you cover a slot with a shingle, add or remove the printed modifiers to your modifier deck.
	Add the chosen classes cards to your card pool.
	Increase your life tracker to reflect the new maximum life.
	
	%Shingles passgenau auf das brett zuschneiden?
	%Dreiecke, Squares, Fünfecke an de seiten?
	%Runde ecken etc?
	
	\section{Events}
	Events will happen to you between each fight.
	Events can have many interesting effects and will warp how you choose to play your game.
	
	
	\section{The Storybook}
	The storybook contains all of the Stories you can play in Wheel of Elements.
	Each story will take about 2 hours to play through, contain 3 Encounters and one text containing a short story for why you are fighting.
	
	\subsection{Stories}
	Each story begins with a plot hook, which should be narrated by a player.
 	The text is separated into four sections. After reading a section, start one of the three fights. After finishing a fight, read another section, draw an event, and level up. After finishing the second fight, repeat this process.
 	%Flowchart? with nice pics?
 	
	
	\subsection{Encounters}
	Each story comes with three enemy encounters. An encounter presents which enemies should be fought and in which arrangement.
	There are always 5 enemy setups described per encounter. Each player gets one of those setups.
	These enemy setups are distributed in clockwise seating order. %Please get  a picture for this
	If you are playing with 
	
	
	\section{Enemies}
	\subsection{Enemy Keywords}
	\section{Combat Round}
	\subsection{Player Turn}
	Shuffle your Action deck and your modifier deck before the combat starts. Empty the Elemental Wheel.\\
	Start of Turn\\
	Draw up to your maximum hand size.
		Empty draw pile - When your draw pile is empty and you need to draw, your discard pile becomes your new draw pile. Shuffle it.
	Roll the Monster Die
	Start of turn and start of combat abilities trigger. Players may decide order.
	Play\\
	Players may play up to 2 cards each turn.
	To play a card:
	Pay its mandatory element cost
	Choose if you want to play the optional costs.
	Resolve the card.
	Put the played card into the discard pile.
	
	\subsection{Enemy Turn}
	Your leftmost enemy, that didnt act yet, uses their action. 
	Resolve the action. 
	Repeat until all your enemies acted, even if you die.
		
	\subsection{Cleanup step}
	Remove all single Turn Boons\&Banes (They are Shield, Player Powerful, Advantage etc).\\
	Apply Poison Damage.\\
	Move your leftmost enemy to the player to your left. This happens simultaneous. You cannot give away an enemy given to you this way.\\
	If you are dead give all your remaining enemies to your left.
	
	\section{Playing a Card}
	\section{Boons \& Banes}
	Boons
	Powerful
	Advantage
	Shield
	Flurry
	
	Banes
	Poison
	Vulnerable
	Disarm
	
	\section{Optional Rules}
	You can start with a Class level 1 and its corresponding class level 2
	
	\section{Background Info for Devs}
	\subsection{Math For nerds}
	\subsection{Card List}
\end{document}